
\subsection{Tipos de grafo}
\label{sec:tipos_de_grafo}

Dependiendo del tipo de grafo, el objeto matemático es ligeramente distinto. Salvo que se diga lo contrario, se presume que todo grafo es simple, no dirigido y no ponderado, de acuerdo con las definiciones a continuación.

\subsubsection{Dirigido o no dirigido}
\label{sec:dirigido_vs_no_dirigido}

Dependiendo de la estructura del conjunto de arcos, un grafo es dirigido o no dirigido:

\begin{itemize}
  \item Si los arcos de un grafo se dirigen de un vértice hacia otro, el grafo es \concept{dirigido}. En ese caso, \(E \subseteq V \times V\) y los arcos son pares \emph{ordenados} de la forma \((u,v)\).
  \item Si los arcos de un grafo no tienen dirección, el grafo es \concept{no dirigido}. Si \(\binom{V}{2}\) denota el conjunto de subconjuntos de cardinalidad \(2\) de \(V\), entonces \(E \subseteq \binom{V}{2}\), es decir, los arcos son pares \emph{no ordenados} de la forma \(\{u,v\}\).
\end{itemize}

Indistintamente de si el grafo es dirigido o no, es común denotar los arcos como \(uv\) cuando no hace falta ser explícito en su estructura, en lugar de \((u,v)\) y \(\{u,v\}\), respectivamente.

\subsubsection{Simple o multigrafo}
\label{sec:simple_o_multigrafo}

Un \concept{lazo} (\textit{loop}) es un arco que conecta un vértice consigo mismo, de la forma \(\{v, v\}\) en un grafo no dirigido o \((v, v)\) en un grafo dirigido.

Un grafo es \concept{simple} si no tiene lazos y además hay a lo sumo un arco entre cada par de vértices. Salvo que explícitamente se diga lo contrario, los conceptos en estos apuntes suponen grafos simples.

Si un grafo puede tener lazos, se dice que es un \concept{grafo simple con lazos}. Su conjunto de arcos es de la forma \(E \subseteq \binom{V}{2} \cup \{\{a\}|a\in V\}\) y los lazos son conjuntos \textit{singleton} con un solo vértice.

Si un grafo puede tener múltiples arcos entre el mismo par de vértices, se dice que es un \concept{multigrafo}. En ese caso, \(E \subseteq \binom{V}{2} \times \mathbb{N}\) es el conjunto de arcos, de la forma \((\{u, v\}, n)\) donde \(u,v\in V\) y \(n \in \mathbb{N}\). El número asociado a cada arco no es un peso, a diferencia de en un \hyperref[sec:grafos_ponderados]{grafo ponderado}. Es un diferenciador, para discernir entre arcos que conectan el mismo par de vértices.

\subsubsection{Grafos ponderados}
\label{sec:grafos_ponderados}

Un grafo \concept{ponderado} es una tripla \(G = (V, E, w)\), donde \((V, E)\) es un grafo (de cualquiera de los tipos anteriormente definidos) y \(w\colon E \to \mathbb{R}\) es una función de pesos que asigna un número real a cada arco.
